% ==============================================================================
% Fichier : front/01_avant_propos.tex : FORN (LPI342)
% ==============================================================================

\chapter*{Avant-propos}
\addcontentsline{toc}{chapter}{Avant-propos}
\markboth{Avant-propos}{Avant-propos}

Ce livre est le support de l'unité d'enseignement \textbf{LPI342, Forensique
Numérique}, du Master Professionnel en Sécurité des Systèmes Informatiques et
des Données de l'Université de Garoua (Faculté des Sciences, Département de
Mathématiques et Informatiques).

\medskip
La forensique numérique est l'investigation à vocation probatoire~: établir ce
qui s'est passé \emph{et} pouvoir le prouver. Là où le hacking éthique apprend
à simuler l'attaque, et la supervision à la détecter, FORN apprend à
\textbf{enquêter}~: acquérir sans altérer, analyser disque, mémoire, réseau,
mobile et nuage, reconstituer un faisceau, et restituer par un rapport et
devant un tribunal. Le fil conducteur est une enquête unique, l'\textbf{opération
TSANAGA}, menée à Garoua par le cabinet \textbf{NYONG SÉCURITÉ} pour le compte
du parquet, sur les supports saisis chez un ancien employé de l'entreprise
fictive \textbf{LiZbiZ}, soupçonné d'avoir revendu le fichier des clients.

\medskip
Chaque notion est ancrée à la fois dans les standards internationaux (ISO/IEC
27037, 27042 et 27043, NIST SP 800-86 et 800-101, RFC 3227, guide ACPO) et
dans le droit camerounais et africain (lois n°~2010/012 et 2024/017,
Conventions de Malabo et de Budapest). Un fil théorique parcourt le livre~: le
\textbf{Trilemme CRO} (Confidentialité, Fiabilité, Opposabilité) et la
\textbf{chaîne de possession}, par lesquels se juge toute trace.

\medskip
Les cas sont fictifs~: les noms, les dates et les montants le sont~; les
mécanismes techniques et les références normatives sont exacts. Le livre
s'accompagne de l'environnement \textbf{NDAP SUKULU}, qui fait pratiquer
chaque chapitre.

\bigskip
\noindent \IfFileExists{../../Logos/signature_avant_propos.png}{\includegraphics[height=1.6cm]{../../Logos/signature_avant_propos.png}}{\textit{MINKA MI NGUIDJOÏ Thierry Emmanuel}}
